\PassOptionsToPackage{dvipsnames,svgnames,table}{xcolor}
\documentclass[11pt]{article}
\usepackage[a4paper,margin=22mm]{geometry}
\usepackage[no-math]{fontspec}
\setmainfont{DejaVu Serif}
\usepackage{amsmath,amssymb,amsthm,mathtools,graphicx,xcolor,hyperref,xurl,xspace,ifthen,enumerate}
\usepackage[shortlabels]{enumitem}
\setlength{\emergencystretch}{4em}
\AtBeginDocument{\special{pdf:mapline pzdr Dingbats <uzdr.pfb}}
\SetMathAlphabet{\mathbf}{normal}{OT1}{cmr}{bx}{n}
\SetMathAlphabet{\mathbf}{bold}{OT1}{cmr}{bx}{n}
\newtheorem{theo}{Theorem}[section]
\newtheorem{lemm}[theo]{Lemma}
\newtheorem{thm}[theo]{Theorem}
\newtheorem{proposition}[theo]{Proposition}
\newtheorem{corollary}[theo]{Corollary}
\newtheorem{definition}[theo]{Definition}
\newtheorem{theorem}[theo]{Theorem}
\newtheorem*{remas*}{Remarks}
\newtheorem{remas}[theo]{Remarks}
\newtheorem*{exam*}{Example}
\newtheorem{ques}[theo]{Question}
\providecommand{\og}{«\nobreakspace}
\providecommand{\fg}{\nobreakspace»}
\newtheorem{lem}[theo]{Lemma}
\newtheorem{lemma}[theo]{Lemma}
\newtheorem{prop}[theo]{Proposition}
\newtheorem{coro}[theo]{Corollary}
\newtheorem{corol}[theo]{Corollary}
\newtheorem*{conj*}{Conjecture}
\newtheorem{conj}[theo]{Conjecture}
\newtheorem{conjecture}[theo]{Conjecture}
\newtheorem{Proposition}[theo]{Proposition}
\newtheorem{theoreme}[theo]{Theorem}
\theoremstyle{definition}
\newtheorem{defi}[theo]{Definition}
\newtheorem{exam}[theo]{Example}
\newtheorem{exem}[theo]{Example}
\newtheorem{remark}[theo]{Remark}
\newtheorem{example}[theo]{Example}
\newtheorem{exams}[theo]{Examples}
\newtheorem{rema}[theo]{Remark}
\newtheorem{nota}[theo]{Notation}
\newtheorem{notation}[theo]{Notation}
\newtheorem*{notas}{Notations}
\newtheorem{result}[theo]{Result}
\newtheorem{question}[theo]{Question}
\newtheorem{prob}[theo]{Problem}
\newtheorem*{rema*}{Remark}
\newtheorem*{defi*}{Definition}
\newtheorem*{theo*}{Theorem}
\newtheorem*{lemm*}{Lemma}
\newtheorem*{prop*}{Proposition}
\newcommand{\enoncename}{}
\newtheorem{corpusenonce}[theo]{\enoncename}
\newenvironment{enonce}[2][]{\renewcommand{\enoncename}{#2}\begin{corpusenonce}}{\end{corpusenonce}}
\newenvironment{enonce*}[2][]{\par\medskip\noindent\textbf{#2.}\itshape}{\par\medskip}
\newenvironment{proofc}[1][\proofname]{\begin{proof}[#1]}{\end{proof}}
\newcommand{\Mc}{\mathcal{M}}
\newcommand{\pointir}{.\ ---\ }
\newcommand{\equalenv}[2]{\expandafter\let\csname#1\expandafter\endcsname\csname#2\endcsname\expandafter\let\csname end#1\expandafter\endcsname\csname end#2\endcsname}
\providecommand{\diagbox}[3][]{\begin{tikzpicture}[x=1em,y=1em,baseline=(current bounding box.center)]\useasboundingbox(0,0) rectangle(3,1.4);\draw(0,1.4)--(3,0);\node at(.5,.3){#2};\node at(2.5,1.1){#3};\end{tikzpicture}}
\providecommand{\eg}{e.g.\ }
\providecommand{\ie}{i.e.\ }
\providecommand{\cf}{cf.\ }
\providecommand{\longthanks}[1]{\par\smallskip\noindent #1\par\smallskip}
\providecommand{\qbinom}[2]{\genfrac{[}{]}{0pt}{}{#1}{#2}}

\providecommand{\mathof}[1]{\mathrm{#1}}

\numberwithin{equation}{section}


\def\endappendix{}


\usepackage{latexsym,amsmath,amssymb,verbatim,tikz,tikz-cd}
\usepackage{mathtools}
\usepackage{graphicx,xcolor}

\usepackage[english]{babel}
\usepackage{amssymb}
\usepackage{amsmath}
\usepackage{amsthm}
\usepackage{graphicx}
\usepackage{csquotes}
\usepackage{bm}
\usepackage{enumitem}
\usepackage{subcaption}
\usepackage{caption}
\renewcommand{\thesubfigure}{\Alph{subfigure}}
\usepackage{dirtytalk}
\providecommand{\nopunct}{\spacefactor1007\relax}
\newtheorem{problem}[theo]{Problem}

\newcommand{\sym}{\mathrm{Sym}}
\newcommand{\qsym}{\mathrm{QSym}}
\newcommand{\Des}{\operatorname{Des}}
\newcommand{\Asc}{\operatorname{Asc}}
\newcommand{\Syt}{\operatorname{SYT}}
\newcommand{\Short}{\operatorname{Short}}
\newcommand{\Loop}{\operatorname{Loop}}
\newcommand{\Stable}{\operatorname{Stable}}
\newcommand{\len}{\operatorname{len}}
\newcommand{\height}{\operatorname{height}}
\newcommand{\core}{\operatorname{core}}
\newcommand{\res}{\operatorname{res}}
\newcommand{\row}{\operatorname{row}}
\newcommand{\insrt}{\operatorname{insert}}
\newcommand{\UP}{\operatorname{UP}}
\newcommand{\DOWN}{\operatorname{DOWN}}


\newcommand{\ZZ}{\mathbb{Z}}
\newcommand{\NN}{\mathbb{N}}
\newcommand{\RR}{\mathbb{R}}
\newcommand{\QQ}{\mathbb{Q}}
\newcommand{\PP}{\mathbb{P}}

\newcommand{\F}{\mathcal{F}}
\newcommand{\Q}{\mathcal{Q}}
\newcommand{\A}{\mathcal{A}}
\newcommand{\M}{\mathcal{M}}
\newcommand{\I}{\mathcal{I}}
\newcommand{\PM}{\mathcal{PM}}
\newcommand{\Pcal}{\mathcal{P}}


\newcommand{\odds}{\mathrm{odds}}

\makeatletter
\newcommand{\proofstep}[1]{
  \par
  \addvspace{\medskipamount}
  \textbf{#1\@addpunct{.}}\enspace\ignorespaces
}

\newcommand\restr[2]{{
  \left.\kern-\nulldelimiterspace 
  #1 
  \vphantom{\big|} 
  \right|_{#2} 
  }}
\makeatother












\title{Inserting a short chord at the same position preserves Knuth equivalence of matchings}
\author{Avichai Marmor}
\date{doi:10.5802/alco.351; CC BY 4.0}
\begin{document}
\maketitle
\noindent\textit{Editorial note.} The counted statement is the author's labelled lemma environment at publisher-source.tex lines 1167-1170; its proof environment opens at line 1172, the line after the blank separator, so the pairing is by adjacency and the retained proof range 1172-1176 is the whole environment.\par\medskip
\noindent\textit{Editorial note.} The lemma lem:insert i equivalent insert j (statement, source lines 1157-1160) is retained as context because the counted proof cites it at source line 1173; it is the only outside label of the retained proof.\par\medskip
\noindent\textit{Editorial note.} This article carries no label-number dump and the retained counter values are left empty for the whole source, so the printed numbering of the card is the wrapper's default; the identity of the counted item is its author label lem:insert same chord to equivalent remains equivalent.\par\medskip
\noindent\textit{Editorial scope.} One original result of the pinned publisher article, counted once: invariance of Knuth equivalence under inserting a short chord at a fixed position (author label lem:insert same chord to equivalent remains equivalent).\par\medskip
\noindent\textit{Allowed standard prerequisites (editorial).}\begin{itemize}
\item Insertion of a short chord. insrt\_i(m) is the matching obtained by inserting a short chord joining i and i+1 and shifting the later vertices by two.. It is the operation applied to both matchings in the statement.
\item Knuth equivalence generated by elementary transformations. Two matchings are Knuth equivalent when one is obtained from the other by a sequence of elementary Knuth transformations.. It is both the hypothesis and the conclusion, and the proof inducts on the length of such a sequence.
\item Independence of the insertion position. Inserting a short chord at any two positions gives Knuth equivalent matchings.. It lets the proof restrict to the last position.
\end{itemize}\medskip
\begin{lemma}
\label{lem:insert i equivalent insert j}
    Let $m \in \M_N$ be a matching and let $1 \le i,j \le N+1$ be indices. Then $\insrt_i(m)$ is Knuth equivalent to $\insrt_j(m)$.
\end{lemma}

\begin{lemma}
\label{lem:insert same chord to equivalent remains equivalent}
    Let $m_1,m_2 \in \M_N$ be Knuth equivalent matchings, and let $1 \le i \le N+1$ be an index. Then $\insrt_i(m_1)$ and $\insrt_i(m_2)$ are Knuth equivalent.
\end{lemma}

\begin{proof}
    By Lemma~\ref{lem:insert i equivalent insert j}, we may assume that $i = N+1$. Assume that $m_2 = \varphi_1 \cdots \varphi_\ell (m_1)$ for some elementary Knuth transformations $\varphi_1, \dots, \varphi_\ell$. We prove the statement by induction on $\ell$. The statement is obvious for $\ell = 0$.

    Next, assume that $\ell > 0$. By the induction hypothesis, we may assume that $\insrt_{N+1}(m_2)$ is equivalent to $\insrt_{N+1}(\varphi_\ell(m_1))$. Therefore, it suffices to show that $\insrt_{N+1}(m_1)$ is equivalent to $\insrt_{N+1}(\varphi_\ell(m_1))$. Assume that $(j,j+1),\; (j+2,j') \in m_1$ for some $1 \le j,j' \le N$, and that $(j,j'),\; (j+1,j+2) \in \varphi_\ell(m_1)$, as other types of elementary Knuth transformations are handled similarly. Notice that $\insrt_{N+1}(m_1)$ and $\insrt_{N+1}(\varphi_\ell(m_1))$ have all but four chords in common. Specifically, $\insrt_{N+1}(m_1)$ contains the chords $(j,j+1)$ and $(j+2,j')$, while $\insrt_{N+1}(\varphi_\ell(m_1))$ contains the chords $(j,j')$ and $(j+1,j+2)$. Therefore, there exists an elementary Knuth transformation $\varphi$, such that $\varphi(\insrt_{N+1}(m_1)) = \insrt_{N+1}(\varphi_\ell(m_1))$. Thus, the matchings $\insrt_{N+1}(m_1)$ and $\insrt_{N+1}(\varphi_\ell(m_1))$ are equivalent, as required.
\end{proof}

\end{document}
